\documentclass[12pt,a4paper]{article}
\usepackage{amsmath,amssymb,amsthm}
\usepackage{hyperref}
\usepackage{geometry}
\geometry{margin=2.5cm}
\usepackage{enumitem}
\usepackage{booktabs}

\newtheorem{theorem}{Theorem}[section]
\newtheorem{lemma}[theorem]{Lemma}
\newtheorem{corollary}[theorem]{Corollary}
\newtheorem{proposition}[theorem]{Proposition}
\theoremstyle{definition}
\newtheorem{definition}[theorem]{Definition}
\theoremstyle{remark}
\newtheorem{remark}[theorem]{Remark}

\title{Primordial Nucleosynthesis from Recognition Science:\\
BBN Yields from the $\varphi$-Ladder and Ledger Bookkeeping}

\author{Jonathan Washburn\\
Recognition Science Research Institute\\
Austin, Texas\\
\texttt{washburn.jonathan@gmail.com}}

\date{March 2026}

\begin{document}
\maketitle

\begin{abstract}
We derive Big Bang Nucleosynthesis (BBN) yields from the Recognition
Science (RS) framework.  The baryon-to-photon ratio
$\eta = n_b/n_\gamma \approx 6.1 \times 10^{-10}$ is connected to
the $\varphi$-ladder: the closest integer-rung match is
$\varphi^{-44} \approx 6.38 \times 10^{-10}$, within $4.5\%$ of the
observed value.  With $\eta$ fixed, the primordial helium-4 mass
fraction $Y_p \approx 0.245$, the deuterium abundance
$D/H \approx 2.5 \times 10^{-5}$, and the helium-3 abundance
${}^3\mathrm{He}/H \approx 10^{-5}$ follow from the standard
nuclear reaction network.  The neutron-to-proton freeze-out ratio
$n/p \approx 1/7$ is set by the neutron--proton mass difference
$Q = 1.293\;\mathrm{MeV}$ and the weak interaction freeze-out
temperature $T_f \approx 0.8\;\mathrm{MeV}$.  The cosmological
lithium-7 problem---where the predicted
${}^7\mathrm{Li}/H \approx 5 \times 10^{-10}$ is
${\sim}\,3\times$ higher than observed---is addressed through
post-BBN stellar depletion mechanisms.  The effective number of neutrino species $N_{\mathrm{eff}} = 3.044$
follows from RS's three-generation theorem~\cite{RS_Axioms}
($f(D) = D = 3$ for $D = 3$ spatial dimensions) combined with
standard neutrino decoupling physics.
\end{abstract}

%% ============================================================
\section{Introduction}
\label{sec:intro}
%% ============================================================

Big Bang Nucleosynthesis~\cite{Fields2020} is one of the three
pillars of the hot Big Bang model, along with the cosmic microwave
background and the Hubble expansion.  BBN predictions depend
sensitively on:
\begin{enumerate}
  \item The baryon-to-photon ratio $\eta$.
  \item The neutron lifetime $\tau_n$.
  \item The number of light neutrino species $N_\nu$.
\end{enumerate}

In the standard model, $\eta$ is a free parameter fitted to
observations.  In Recognition Science, $\eta$ is connected to the
$\varphi$-ladder through the CP-violating phase structure of the
8-tick epoch.  This paper derives the BBN yields from that
connection.

\subsection{Recognition Science Framework}
\label{sec:framework}

The BBN derivation rests on the unique cost functional forced by
three axioms~\cite{RS_Axioms}.

\begin{definition}[$J$-Cost Functional]
For any positive ratio $x > 0$:
\begin{equation}
  J(x) = \tfrac{1}{2}(x + x^{-1}) - 1.
\end{equation}
\end{definition}

\noindent\textbf{Axiom A1 (Normalization):} $J(1) = 0$.\\
\textbf{Axiom A2 (Recognition Composition Law, RCL):}
$J(xy)+J(x/y)=2J(x)J(y)+2J(x)+2J(y)$ for all $x,y>0$.\\
\textbf{Axiom A3 (Calibration):} $J''_{\log}(0) = 1$.

\begin{theorem}[Cost Uniqueness]
\label{thm:cost-unique}
The continuous solution to \textup{(A1)--(A3)} is unique:
$J(x) = \tfrac{1}{2}(x + x^{-1}) - 1$.
\end{theorem}

\begin{proof}[Proof sketch]
Setting $H(t)=J_{\log}(t)+1$ transforms A2 into the d'Alembert
equation $H(s+t)+H(s-t)=2H(s)H(t)$.  The Acz\'el classification
theorem~\cite{Aczel1966} uniquely selects $H(t)=\cosh t$, giving
$J(x)=\tfrac{1}{2}(x+x^{-1})-1$.
\end{proof}

\noindent Dimension $D=3$ is forced by the 8-tick recognition epoch
($2^D=8$).  This determines the number of fermion generations:
$N_{\mathrm{gen}} = f(D) = D = 3$, giving exactly three light neutrino
species and an irreducible CKM CP-violating phase---both essential
inputs to BBN.

%% ============================================================
\section{The Baryon-to-Photon Ratio}
\label{sec:eta}
%% ============================================================

\begin{proposition}[$\eta$ and the $\varphi$-Ladder]
\label{prop:eta}
The baryon-to-photon ratio lies near rung $-44$ on the
$\varphi$-ladder:
\begin{equation}
  \boxed{\varphi^{-44} \approx 6.38 \times 10^{-10}
  \approx \eta_{\mathrm{obs}}.}
\end{equation}
\end{proposition}

\begin{proof}[Structural observation]
The proximity is established by computing the $\varphi$-rung
of the observed $\eta$:
\begin{equation}
  n = \frac{\ln(1/\eta_{\mathrm{obs}})}{\ln\varphi}
  = \frac{\ln(1.639 \times 10^{9})}{0.4812}
  \approx 44.1.
\end{equation}
The nearest integer rung is $n = 44$, giving
$\varphi^{-44} \approx 6.38 \times 10^{-10}$.  This is a
\emph{structural observation} that the observed $\eta$ is close
to a $\varphi$-ladder rung, not a derivation of $\eta$ from first
principles.  A genuine RS derivation would compute the
baryogenesis rate directly from the ledger's CP-violating
structure; see the companion EWPT paper for a structural estimate.
\end{proof}

\begin{remark}[Observational Comparison]
The observed value from Planck CMB analysis is
$\eta_{\mathrm{obs}} = (6.104 \pm 0.058) \times
10^{-10}$~\cite{Planck2020}.  The RS prediction
$\varphi^{-44} \approx 6.38 \times 10^{-10}$ is $4.5\%$ above
the central value, a ${\sim}\,4.7\sigma$ tension.  This suggests
that the true $\varphi$-ladder connection may involve a
non-integer rung or a multiplicative correction factor from the
detailed baryogenesis mechanism.  The derivation of the exact
exponent from first principles remains an active area of RS
development.
\end{remark}

\begin{definition}[Baryon Density Parameter]
\begin{equation}
  \Omega_b h^2 = 3.65 \times 10^7 \cdot \eta \approx 0.0233,
\end{equation}
where $h = H_0/(100\;\mathrm{km/s/Mpc})$.  The observed value is
$\Omega_b h^2 = 0.02237 \pm 0.00015$~\cite{Planck2020}.
\end{definition}

%% ============================================================
\section{Neutron-to-Proton Freeze-Out}
\label{sec:np}
%% ============================================================

\begin{definition}[Neutron--Proton Mass Difference]
$Q = m_n - m_p = 1.293\;\mathrm{MeV}$.
\end{definition}

\begin{theorem}[Freeze-Out Ratio]
\label{thm:np}
At the weak interaction freeze-out temperature
$T_f \approx 0.8\;\mathrm{MeV}$:
\begin{equation}
  \frac{n}{p}\bigg|_{\mathrm{freeze}} = e^{-Q/T_f}
  \approx e^{-1.293/0.8} \approx 0.20 \approx \frac{1}{5}.
\end{equation}
After neutron decay during the interval between freeze-out and
deuterium formation ($t \approx 180\;\mathrm{s}$):
\begin{equation}
  \frac{n}{p}\bigg|_{\mathrm{BBN}} \approx
  \frac{(1/5) \cdot e^{-180/880}}{1 + (1/5)(1 - e^{-180/880})}
  \approx \frac{1}{7}.
\end{equation}
\end{theorem}

\begin{proof}
Before freeze-out, the $n \leftrightarrow p$ reactions (mediated by
weak interactions) maintain equilibrium: $n/p = e^{-Q/T}$.  At
$T = T_f$, the weak interaction rate drops below the Hubble rate,
freezing the ratio at ${\sim}\,1/5$.  During the ${\sim}\,180$~s
before deuterium synthesis begins, neutrons decay with lifetime
$\tau_n \approx 880\;\mathrm{s}$: the survival fraction is
$e^{-180/880} \approx 0.815$, reducing $n/p$ from $1/5$ to
approximately $1/7$.
\end{proof}

\begin{remark}
The freeze-out temperature $T_f$ depends on the Fermi constant
$G_F$, which in RS is derived from the $\varphi$-ladder (see the
companion paper on electroweak parameters).  The neutron--proton
mass difference $Q$ depends on $m_d - m_u$ and the electromagnetic
self-energy, both of which are $\varphi$-ladder quantities.
\end{remark}

%% ============================================================
\section{Primordial Abundances}
\label{sec:abundances}
%% ============================================================

\begin{theorem}[Helium-4 Mass Fraction]
\label{thm:helium}
\begin{equation}
  Y_p = \frac{2(n/p)}{1 + n/p} = \frac{2 \cdot 1/7}{1 + 1/7}
  = \frac{2/7}{8/7} = \frac{1}{4} = 0.250.
\end{equation}
Including corrections for incomplete burning:
$Y_p \approx 0.245$.
\end{theorem}

\begin{proof}
Almost all neutrons at the time of BBN are incorporated into
helium-4 (the most tightly bound light nucleus).  Each helium-4
nucleus contains 2 neutrons and 2 protons.  With $n/p = 1/7$,
for every 2 neutrons there are 14 protons; the 2 neutrons pair
with 2 protons to form 1 helium-4 nucleus (mass 4), leaving 12
protons (mass 12).  The helium mass fraction is $4/(4 + 12) = 1/4$.
\end{proof}

\begin{theorem}[Light Element Abundances]
For $\eta \approx 6.1 \times 10^{-10}$, the standard BBN reaction
network yields:
\begin{align}
  Y_p &\approx 0.245, \\
  D/H &\approx 2.5 \times 10^{-5}, \\
  {}^3\mathrm{He}/H &\approx 1.0 \times 10^{-5}, \\
  {}^7\mathrm{Li}/H &\approx 5.0 \times 10^{-10}.
\end{align}
\end{theorem}

\begin{remark}[Observational Comparison]~

\medskip
\begin{center}
\renewcommand{\arraystretch}{1.3}
\begin{tabular}{@{}lccc@{}}
\toprule
\textbf{Species} & \textbf{RS/BBN prediction} &
\textbf{Observed} & \textbf{Agreement?} \\
\midrule
${}^4$He ($Y_p$) & 0.245 & $0.2449 \pm 0.0040$ & Yes \\
D/H & $2.5 \times 10^{-5}$ & $(2.527 \pm 0.030) \times 10^{-5}$
& Yes \\
${}^3$He/H & $1.0 \times 10^{-5}$ & $(1.1 \pm 0.2) \times 10^{-5}$
& Yes \\
${}^7$Li/H & $5.0 \times 10^{-10}$ &
$(1.6 \pm 0.3) \times 10^{-10}$ & Discrepant \\
\bottomrule
\end{tabular}
\end{center}
\end{remark}

%% ============================================================
\section{The Lithium Problem}
\label{sec:lithium}
%% ============================================================

\begin{remark}[Lithium Discrepancy]
The predicted ${}^7\mathrm{Li}/H \approx 5 \times 10^{-10}$ is
${\sim}\,3\times$ the observed Spite-plateau value
$1.6 \times 10^{-10}$~\cite{Spite1982}.  This is an observational
discrepancy, not an RS prediction; it is common to the standard BBN
framework.
\end{remark}

\begin{remark}[RS Perspective]
The discrepancy is attributed to post-BBN astrophysical depletion,
not to errors in the primordial prediction.  Lithium-7 is fragile
($B/A \approx 5.6\;\mathrm{MeV}$) and is destroyed by stellar
burning at temperatures $T > 2.5 \times 10^6\;\mathrm{K}$.
Pop~II halo stars, where the primordial lithium is measured, have
convective zones that can mix surface lithium into hotter interior
layers.  Depletion factors of ${\sim}\,3\times$ over
${\sim}\,10\;\mathrm{Gyr}$ are consistent with detailed stellar
models~\cite{Korn2006}.  RS does not modify the nuclear reaction
network; the primordial abundance remains
${}^7\mathrm{Li}/H \approx 5 \times 10^{-10}$.
\end{remark}

%% ============================================================
\section{Effective Neutrino Species}
\label{sec:neff}
%% ============================================================

\begin{theorem}[Neutrino Species from BBN]
$N_{\mathrm{eff}} = 3.044$.
\end{theorem}

\begin{proof}
RS forces exactly 3 fermion generations from the face-pairing
structure of the $D = 3$ lattice: the number of face pairs of a
$D$-cube is $D$, so $N_{\mathrm{gen}} = 3$.  The effective number
$N_{\mathrm{eff}}$ exceeds 3 slightly because neutrinos are not
fully decoupled at $e^+e^-$ annihilation; some of the $e^+e^-$
energy heats the neutrinos.  The standard computation
(incorporating finite-temperature QED corrections and neutrino
oscillations) gives $N_{\mathrm{eff}} = 3.044$~\cite{Mangano2005}.
RS does not modify this calculation, since the ledger's discrete
structure operates at the Planck scale, far above the MeV-scale
physics of neutrino decoupling.
\end{proof}

\begin{remark}
Planck measures $N_{\mathrm{eff}} = 2.99 \pm 0.17$.  The RS/SM
prediction of 3.044 is consistent within $0.3\sigma$.
\end{remark}

%% ============================================================
\section{The BBN Reaction Network}
\label{sec:network}
%% ============================================================

The key reactions in primordial nucleosynthesis:
\begin{align}
  n + p &\to D + \gamma, \label{eq:r1} \\
  D + D &\to {}^3\mathrm{He} + n, \label{eq:r2} \\
  D + D &\to T + p, \label{eq:r3} \\
  {}^3\mathrm{He} + D &\to {}^4\mathrm{He} + p, \label{eq:r4} \\
  T + D &\to {}^4\mathrm{He} + n, \label{eq:r5} \\
  {}^3\mathrm{He} + {}^4\mathrm{He} &\to {}^7\mathrm{Be} + \gamma.
  \label{eq:r6}
\end{align}
${}^7\mathrm{Be}$ subsequently captures an electron to produce
${}^7\mathrm{Li}$.

\begin{remark}
The cross-sections for these reactions are determined by nuclear
physics.  RS does not modify the cross-sections---the RS contribution
is to fix the initial conditions ($\eta$, $n/p$, $N_\nu$) that
determine the yields.
\end{remark}

%% ============================================================
\section{Predictions}
\label{sec:predictions}
%% ============================================================

\begin{enumerate}
  \item \textbf{$\eta \approx \varphi^{-44}$}: A near-integer-rung
  prediction on the $\varphi$-ladder.  The $4.5\%$ residual between
  $\varphi^{-44}$ and the observed $\eta$ is a target for
  first-principles derivation of the exact exponent.

  \item \textbf{No sterile neutrinos in BBN}: $N_{\mathrm{eff}}$
  should remain consistent with 3.044.  Light sterile neutrinos would
  increase $N_{\mathrm{eff}}$ and $Y_p$.

  \item \textbf{Lithium depletion}: The primordial
  ${}^7\mathrm{Li}/H$ should be recoverable from pristine gas clouds
  (e.g., damped Lyman-$\alpha$ systems) at
  ${\sim}\,5 \times 10^{-10}$, not the depleted stellar value
  of $1.6 \times 10^{-10}$.
\end{enumerate}

%% ============================================================
\section{Conclusion}
\label{sec:conclusion}
%% ============================================================

Primordial nucleosynthesis in Recognition Science reproduces all
standard BBN yields with the baryon-to-photon ratio connected to
the $\varphi$-ladder at $\eta \approx \varphi^{-44} \approx
6.38 \times 10^{-10}$, within $4.5\%$ of the observed value.
The helium-4, deuterium, and helium-3 abundances match
observations.  The lithium-7 discrepancy is attributed to
post-BBN stellar depletion.  The number of effective neutrino
species $N_{\mathrm{eff}} = 3.044$ follows from RS's
three-generation theorem combined with standard decoupling physics.
The first-principles derivation of the exact $\varphi$-ladder
exponent for $\eta$ remains an active area of RS development.

%% ============================================================
\begin{thebibliography}{99}

\bibitem{RS_Axioms}
J.~Washburn, ``The Algebra of Reality: A Recognition Science Derivation
of Physical Law,'' \emph{Axioms} \textbf{15}(2), 90 (2026).

\bibitem{Aczel1966}
J.~Acz\'el, \textit{Lectures on Functional Equations and Their Applications},
Academic Press, New York (1966).

\bibitem{Fields2020}
B.~D.~Fields, K.~A.~Olive, T.-H.~Yeh, and C.~Young, ``Big-Bang
Nucleosynthesis after Planck,'' JCAP \textbf{03}, 010 (2020).

\bibitem{Planck2020}
Planck Collaboration, ``Planck 2018 Results.\ VI.\ Cosmological
Parameters,'' Astron.\ Astrophys.\ \textbf{641}, A6 (2020).

\bibitem{Spite1982}
F.~Spite and M.~Spite, ``Abundance of Lithium in Unevolved Halo
Stars and Old Disk Stars: Interpretation and Consequences,''
Astron.\ Astrophys.\ \textbf{115}, 357 (1982).

\bibitem{Korn2006}
A.~J.~Korn et~al., ``A Probable Stellar Solution to the
Cosmological Lithium Discrepancy,'' Nature \textbf{442}, 657
(2006).

\bibitem{Mangano2005}
G.~Mangano et~al., ``Relic Neutrino Decoupling Including Flavour
Oscillations,'' Nucl.\ Phys.\ B \textbf{729}, 221 (2005).

\end{thebibliography}

\end{document}
